% =============================================================================
% Lecture 19: Agent Economics
% Agentic Machine Learning Systems (Volume III), Chapter 19
% =============================================================================
% Thesis: the account that matters charges every attempt, failed and canceled
% ones included, to the tasks a contract accepted, and each lever pays only if
% it lowers that ratio or meets a deadline the objective values.
% Hook and synthesis: the token-price fallacy (NM 19.1). Running case: the
% repair agent, shown once as real public repair-agent runs (SWE-bench).
%
% Build: make deck V=3 D=19_agent_economics   (figures: images/make_figs.py)
% Plan:  lecture_map.yml in this folder.
% Numbers: images/calcs.py prints every derived value; sources named there.
% =============================================================================

\documentclass[aspectratio=169, 12pt]{beamer}
\usepackage{../../assets/beamerthememlsys}
\mlsyssetup{
  volume={Volume III},
  chapter={Chapter 19},
  chaptertitle={Agent Economics},
  layer=scale,
  logo={../../assets/img/logo-mlsysbook.png},
  instlogo={../../assets/img/logo-harvard.png},
}
\usepackage[T1]{fontenc}
\usepackage[scaled=0.9]{helvet}
\usepackage{courier}
% Code and real tool output: Inconsolata is narrower than Courier.
\usepackage[varqu,scaled=0.95]{zi4}
\renewcommand{\familydefault}{\sfdefault}
\usepackage{booktabs,amsmath}
\usetikzlibrary{arrows.meta,positioning}
\graphicspath{{images/}}
\setcounter{mlsystotalsections}{5}

\title{Agent Economics}
\author{Vijay Janapa Reddi}
\institute{Harvard University}
\date{}

% A large answer choice for predict and concept-question slides
\newcommand{\choice}[2]{\textbf{\textcolor{mlsysaccent}{#1}}\enspace #2}
% A choice that is highlighted from overlay #1 on (the answer)
\newcommand{\answerat}[3]{\alt<#1->{\colorbox{mlsysaccent!15}{\choice{#2}{#3}}}{\choice{#2}{#3}}}
% A goal or takeaway line with its thumbnail
\newcommand{\thumbline}[2]{%
  \begin{minipage}[c]{2.6cm}\centering\includegraphics[width=2.5cm,height=0.95cm,keepaspectratio]{#1}\end{minipage}%
  \hspace{0.3cm}\begin{minipage}[c]{0.74\textwidth}\small\raggedright #2\end{minipage}}
% One line of a worked example: shown from overlay #1 on
\newcommand{\wstep}[2]{\uncover<#1->{#2}\\[0.2cm]}
% An exact excerpt of real tool output, a config, a document, or source code
\newenvironment{logblock}[1][7.2]{\par\ttfamily\fontsize{#1}{8.8}\selectfont\raggedright\setlength{\parskip}{0pt}}{\par}
% Code indentation of n spaces in the current monospace font
\newcommand{\ind}[1]{\hspace*{#1\fontcharwd\font`0}}
% Cost per accepted task and acceptance rate, as the chapter writes them
\newcommand{\Ceff}{C_{\text{eff}}}
% Colors used in the figures, for text beside them
\definecolor{figviolet}{HTML}{7C3AED}
\definecolor{figred}{HTML}{DC2626}
\definecolor{figgreen}{HTML}{16A34A}
\definecolor{figsky}{HTML}{0284C7}
\definecolor{figamber}{HTML}{D97706}

\begin{document}

% =============================================================================
% TITLE AND HOOK
% =============================================================================
\mlsysnote{[0.5 min] [CORE] Read the thesis aloud once. Lecture 18 asked whether splitting a task across agents pays. Today asks the same question of every lever the runtime controls, with one account. Notes total about 70.5 minutes; the core lecture is about 59.5 minutes. FLEX: notes tagged [OPTIONAL] can be cut. Cut first the price-trend plot, the coverage plot, speculative decoding, the long-horizon benchmark plot, the rate-limit excerpt, the serving measurement, the break-even plot, and the open questions. Every concept the chapter owns stays on a core slide.}
\mlsystitle{Agent Economics}{The account that matters charges every attempt, failed and canceled ones included, to the tasks a contract accepted, and a lever pays only if it lowers that ratio or meets a deadline the objective values.}{cover_tokenomics_blueprint.png}

\begin{frame}{Predict: which agent costs less per repair?}
  \mlsysactivity{Predict}{3 min}
  \begin{columns}[c]
    \begin{column}{0.52\textwidth}
      \scriptsize
      \begin{tabular}{@{}lrr@{}}
        \toprule
        100 tickets, 1 attempt each & \textbf{Frontier} & \textbf{Low price}\\
        \midrule
        \$ per M tokens, in / out & 3 / 15 & 0.30 / 1.20\\
        Turns per attempt & 4 & 7\\
        Tokens per turn, in / out & 16,000 / 800 & 28,000 / 1,200\\
        Environment per turn & \$0.03 & \$0.03\\
        Assessment per attempt & \$0.05 & \$0.05\\
        \textbf{Accepted tickets} & \textbf{75} & \textbf{25}\\
        \bottomrule
      \end{tabular}\\[0.15cm]
      {\small The low-price tokens cost \textbf{a tenth} as much.}\\[0.1cm]
      {\normalsize\choice{A}{Low price is cheaper}\\[0.06cm]\choice{B}{About equal}\\[0.06cm]\answerat{2}{C}{Frontier is cheaper}}
    \end{column}
    \begin{column}{0.45\textwidth}
      \centering
      \mlsyssteps[width=\textwidth,height=4.7cm,keepaspectratio]{hook_step}{2}
    \end{column}
  \end{columns}
  \mlsyscredit{Agent Economics, Napkin Math 19.1 (hypothetical configurations). Per-accepted values: images/calcs.py.}
  \note{[3.5 min] [CORE] Two configurations of the repair agent run the same 100 tickets under the same acceptance contract, one bounded attempt each. Read the table: the low-price model's tokens cost a tenth as much, but it takes 7 turns instead of 4, reads more context per turn, and passes 25 tickets instead of 75. ENGAGE: vote with fingers, one for A, two for B, three for C. Step 2: C. The frontier configuration costs \$0.55 per accepted repair; the low-price one costs \$1.32, 2.4 times as much. Students who chose A compared invoices: the low-price token charges are 71.3 percent lower. WARN: the per-attempt cost falls by only 19.8 percent, because every extra turn pays the environment charge, and acceptance falls by 66.7 percent. Add \$10 of engineer triage per rejected attempt and the gap widens to \$3.88 against \$31.32, 8.1 times. LINK: Lecture 1 defined cost per accepted task for planning; today makes it an account a team can measure, and asks which lever lowers it. The synthesis returns to this table.}
\end{frame}

\begin{frame}{The price of a fixed capability falls 9$\times$ to 860$\times$ a year}
  \centering
  \safeimg[width=0.92\textwidth,height=4.6cm,keepaspectratio]{price_trend.pdf}\\[0.1cm]
  {\small Each line holds a benchmark score fixed, so its fall is price alone.}
  \mlsyscredit{Epoch AI (2025), ``LLM inference prices have fallen rapidly but unequally across tasks'' (lowest price per M tokens, 3:1 input/output blend), via MLSysIM. Fitted rates: images/calcs.py.}
  \note{[1.5 min] [OPTIONAL] Real data. Epoch AI tracks the cheapest model that reaches a fixed score on a benchmark. The gray line is the price of GPT-3-level knowledge on MMLU, 12 times cheaper each year; GPT-4-level knowledge fell 16 times a year and GPT-4-level science 42 times. Across all 21 thresholds in the source, the fitted decline runs from 9 times to about 860 times a year, median 46 times. NARRATE the point: offers like the hook's recur every few months. ENGAGE: if token prices keep falling this fast, which terms of a repair attempt will dominate in two years? (Expected: sandbox holding, tests, review, and the acceptance rate in the denominator.) WARN: the prices are list prices of the cheapest qualifying model, a blend of input and output; they say nothing about turns or acceptance. FLEX: [OPTIONAL]; the hook already makes the point with one pair.}
\end{frame}

\begin{frame}{Recall: when does delegation pay on cost?}
  \mlsysactivity{Retrieval}{1.5 min}
  \begin{columns}[c]
    \begin{column}{0.55\textwidth}
      {\Large\bfseries Close your notes.}\\[0.2cm]
      {\normalsize A delegated tree uses 15$\times$ the tokens of a chat; one agent uses 4$\times$. When can delegation lower cost per accepted task?}\\[0.3cm]
      \uncover<2->{\small Token multiple $15/4 = \mathbf{3.75}$. Delegation must raise acceptance 3.75-fold, so it can break even only where one agent succeeds at most \textbf{\textcolor{mlsysaccent}{26.7\%}} of the time.}
    \end{column}
    \begin{column}{0.42\textwidth}
      \centering\safeimg[width=0.6\textwidth,height=1.1cm,keepaspectratio]{thumb_recall.pdf}\\[0.15cm]
      \begin{logblock}[7.6]
        In our data, agents typically use about 4$\times$ more tokens than chat interactions, and multi-agent systems use about 15$\times$ more tokens than chats.
      \end{logblock}
    \end{column}
  \end{columns}
  \mlsyscredit{Lecture 18 (Napkin Math 18.6). Quote: Hadfield et al., ``How we built our multi-agent research system,'' Anthropic Engineering, 2025 (anthropic.com/engineering/multi-agent-research-system).}
  \note{[1.5 min] [CORE] Retrieval from Lecture 18. Give 45 seconds. The quote is the measurement behind the numbers, from one reported system. Step 2: the delegated cost per accepted task over the single agent's is the token multiple times the single agent's acceptance over the delegated acceptance. With a multiple of 3.75, break-even needs delegated acceptance 3.75 times higher, which is possible only if the single agent's acceptance is at most 26.7 percent. ENGAGE: what did that calculation assume about everything except tokens? (Expected: that cost is proportional to tokens; sandboxes, checks, and reviewers were left out.) Today puts them back in. FLEX: [CORE].}
\end{frame}

\begin{frame}{Why can a cheaper model cost more per finished task?}
  \begin{columns}[c]
    \begin{column}{0.80\textwidth}
      \thumbline{thumb_swe_step2.pdf}{Measure cohort cost per accepted task, failed and canceled work included.}\\[0.12cm]
      \thumbline{thumb_payback.pdf}{Price prefix caching, compaction, and retries under one tariff.}\\[0.12cm]
      \thumbline{thumb_waterfall_step4.pdf}{Bound a speedup by the critical path, and price a model cascade.}\\[0.12cm]
      \thumbline{thumb_kingman_step2.pdf}{Size capacity for trajectory goodput, and price queueing.}\\[0.12cm]
      \thumbline{thumb_budget_step5.pdf}{Conserve a budget across children, and price whole architectures.}
    \end{column}
    \begin{column}{0.17\textwidth}
      \centering\safeimg[width=\textwidth,height=1.9cm,keepaspectratio]{hsa_locator_hsa.pdf}\\[0.05cm]
      {\scriptsize H·S·A: horizon, state, authority all lit}
    \end{column}
  \end{columns}
  \mlsyscredit{Locator: Agent Economics, \S19.1 (book figure).}
  \note{[1.5 min] [CORE] The driving question is the slide title. Five goals, one per segment; each thumbnail is the shape of the figure that teaches it. NARRATE the locator: economics touches every exposure at once. A longer horizon multiplies turns and holding time, carried state is what prefix caching and compaction price, and authority decides which actions need human review, the most expensive term of all. The running case is the repair agent, and once in the first segment we look at real public repair-agent runs and their recorded costs. FLEX: [CORE].}
\end{frame}

% =============================================================================
% SECTION: THE ACCOUNT
% =============================================================================
\mlsysnote{[0.5 min] [CORE] First, what goes in the account, and what divides it.}
\mlsysdivider{What does one accepted task really cost?}{Agent Economics, \S19.2}
\section{Account}

\begin{frame}{An attempt's bill has terms that a token price never touches}
  \centering
  \mlsyssteps[width=0.92\textwidth,height=4.0cm,keepaspectratio]{bill_step}{2}\\[0.1cm]
  {\small $C_{\text{attempt}} = \sum_{t}\big(\text{input} + \text{output} + \text{tools} + \text{sandbox}\big) + \text{verify} + \text{cleanup} + \text{human} + \text{runtime}$}
  \mlsyscredit{Agent Economics, \S19.2.1 (whole-trajectory cost) and Napkin Math 19.1 (hypothetical). Bar values: images/calcs.py.}
  \note{[2.5 min] [CORE] The hook's attempts, one bar each. Step 1: the model tokens alone: \$0.240 for the frontier attempt, \$0.069 for the low-price one, the 71.3 percent saving the invoice shows. Step 2: the rest of the attempt. Each turn holds a sandbox and runs tools at \$0.03, so 7 turns cost \$0.210 against \$0.120, and every attempt pays \$0.05 of assessment. The attempts cost \$0.410 and \$0.329, only 19.8 percent apart. The equation is the chapter's whole-trajectory cost: every charge lands in exactly one category, so a test billed as sandbox time is not counted again as verification. ENGAGE: which terms keep running after the model emits its last token? (Expected: cleanup, verification, and human triage.) WARN: closing the account at the last token is the common error; trailing containers and triage stay on the bill. FLEX: [CORE].}
\end{frame}

\begin{frame}{A real repair run spent \$5.10 per resolved task}
  \centering
  \mlsyssteps[width=0.92\textwidth,height=4.9cm,keepaspectratio]{swe_step}{2}\\[0.02cm]
  {\small Run A: 500 tasks cost \$856.86, and 168 resolved. Unresolved runs averaged \$2.00.}
  \mlsyscredit{SWE-bench Verified submissions (github.com/SWE-bench/experiments): A SWE-agent 0.x; B SWE-agent 1.1, 100-task sample; C, D mini-SWE-agent 2.x. Recorded model-token cost, via MLSysIM.}
  \note{[3 min] [CORE] Real data, and the running case as a real system: public repair-agent runs on SWE-bench Verified, with the cost each agent recorded for every trajectory. These costs are model tokens only; no sandbox or human time. Step 1: the gray bars average the cost of runs that resolved their task: \$1.16 for run A in 2024. A dashboard that reports this number reports survivors. Step 2: the violet bars divide the whole cohort's spend by the tasks resolved. Run A spent \$856.86 on 500 tasks and resolved 168: \$5.10 per resolved task, 4.4 times the survivor average. Unresolved runs cost more, \$2.00 on average, because they ran to the limit. By 2026, runs C and D resolve about three quarters of the tasks, and the gap shrinks to 1.6 times. ENGAGE: why did the gap shrink? (Expected: higher acceptance means fewer failures for each success to carry.) WARN: these runs used different models and scaffolds; compare them only as cohorts, not as a ranking. LINK: Lecture 13 fixed the acceptance contract and matched budgets that make such a comparison fair. FLEX: [CORE].}
\end{frame}

\begin{frame}{Each accepted task carries $(1-\alpha)/\alpha$ rejected ones}
  \begin{columns}[c]
    \begin{column}{0.50\textwidth}
      \centering\safeimg[width=\textwidth,height=5.3cm,keepaspectratio]{multiplier.pdf}
    \end{column}
    \begin{column}{0.47\textwidth}
      \small
      $\Ceff = \dfrac{\sum_i C_{\text{task},i}}{N_{\text{acc}}} = \bar C_{\text{acc}} + \dfrac{1-\alpha}{\alpha}\,\bar C_{\text{rej}}$\\[0.3cm]
      {\footnotesize $\alpha$: share of initiated tasks accepted under one contract by a cutoff. A retry on one ticket adds no accepted task.}\\[0.3cm]
      \begin{mlsyscard}{mlsysaccent}\small A token cut shrinks one term. Acceptance divides \textbf{every} term.\end{mlsyscard}
    \end{column}
  \end{columns}
  \mlsyscredit{Agent Economics, \S19.2.2 (cohort cost per accepted task) and its rejected-per-accepted figure, redrawn at slide size.}
  \note{[2.5 min] [CORE] The cohort form of cost per accepted task. The numerator is every initiated task's account, failed, timed-out, and canceled ones included; the denominator counts distinct tasks accepted under one declared contract by a stated cutoff. Split the numerator by outcome and it becomes the mean accepted cost plus the rejection multiplier times the mean rejected cost. The multiplier is an accounting identity, not a model. Point to the curve: at 25 of 100 accepted each success carries 3 failures; at 75 of 100 it carries 0.33. Below about 30 percent the curve climbs steeply. ENGAGE: run A resolved 168 of 500. How many unresolved runs did each resolved one carry? (Expected: 332 over 168, about 2.) WARN: three successful retries on one ticket still deliver one accepted task. LINK: Lecture 1 owns the per-attempt planning form, expected attempt cost over acceptance probability; this cohort form is what an account measures. FLEX: [CORE].}
\end{frame}

\begin{frame}{True or false: what belongs in the account?}
  \mlsysactivity{Concept question}{2.5 min}
  \begin{columns}[c]
    \begin{column}{0.62\textwidth}
      \small
      \textbf{1.} A child agent is canceled before its bill settles. Its charges leave the numerator.\\[0.1cm]
      \uncover<2->{\footnotesize\textcolor{figred}{\textbf{False.}} Incurred charges stay; money reserved for an unsettled bill is a liability.}\\[0.25cm]
      \textbf{2.} Ten tickets still await test results at the cutoff. The ratio is provisional.\\[0.1cm]
      \uncover<2->{\footnotesize\textcolor{figgreen}{\textbf{True.}} Report the cutoff, the pending count, and the open liability.}\\[0.25cm]
      \textbf{3.} Zero tasks were accepted, so report a cost per accepted task of zero.\\[0.1cm]
      \uncover<2->{\footnotesize\textcolor{figred}{\textbf{False.}} The ratio is undefined; report total spend and the zero.}
    \end{column}
    \begin{column}{0.35\textwidth}
      \begin{mlsyscard}{mlsysaccent}\small\textbf{Report with every $\Ceff$}\\[0.05cm]
      \footnotesize total spend, accepted count with its interval, cutoff, pending tasks, open liabilities\end{mlsyscard}
    \end{column}
  \end{columns}
  \mlsyscredit{Agent Economics, \S19.2.2 and Systems Perspective 19.1.}
  \note{[2.5 min] [CORE] Peer instruction: vote on each, discuss with a neighbor for one minute, vote again. Step 2 reveals. (1) False: a canceled branch's costs remain logged against the task, and any reservation backing work still in flight stays open as liability; the settled total is only a lower bound until it clears. (2) True: an unexecuted test suite is a pending task, not a failure, and a late billing adjustment can move the ratio hours later; the report flags the cutoff and calls the figure provisional. (3) False: with no accepted tasks the ratio is undefined, and a transparent report shows the spend and the zero rather than an artificial floor. WARN: students often treat the ratio as a mean of task costs; it is a cohort ratio, and in small cohorts its variability needs the confidence interval on the accepted count. Land the segment: acceptance divides every charge, so the account counts every attempt and keeps pending work visible. FLEX: [CORE].}
\end{frame}

% =============================================================================
% SECTION: REUSE AND RETRIES
% =============================================================================
\mlsysnote{[0.5 min] [CORE] The account names its terms. The runtime acts most directly on two of them: the inference charge, through reuse, and the number of attempts.}
\mlsysdivider{What do reuse, compaction, and retries buy?}{Agent Economics, \S19.3}
\section{Reuse}

\begin{frame}{A published tariff prices a cached read at a tenth of fresh input}
  \begin{columns}[c]
    \begin{column}{0.64\textwidth}
      {\footnotesize\textbf{Model pricing}, per million tokens (excerpt)}\\[0.1cm]
      {\ttfamily\fontsize{7.4}{9}\selectfont
      \begin{tabular}{@{}lrrrr@{}}
        \toprule
        Model & Base input & 5m cache & Cache hits & Output\\
         & tokens & writes & and refreshes & tokens\\
        \midrule
        Claude Sonnet 4.5 & \$3 / MTok & \$3.75 / MTok & \textbf{\$0.30 / MTok} & \$15 / MTok\\
        Claude Sonnet 5.5 & \$2 / MTok & \$2.50 / MTok & \textbf{\$0.10 / MTok} & \$10 / MTok\\
        \bottomrule
      \end{tabular}}\\[0.2cm]
      \begin{logblock}[7]
        2 Cache hits and refreshes on Claude Opus 5.5 and Claude Sonnet 5.5 are priced at 0.05x the base input price.
      \end{logblock}
    \end{column}
    \begin{column}{0.33\textwidth}
      \small The first row is the chapter's tariff: \$3 fresh, \$0.30 cached, \$15 output.\\[0.2cm]
      The newer row discounts cached reads to \textbf{5\%}.\\[0.2cm]
      {\footnotesize Cache hit: a resent prompt prefix the server still holds.}
    \end{column}
  \end{columns}
  \mlsyscredit{Anthropic, ``Pricing,'' platform.claude.com/docs/en/about-claude/pricing, retrieved 2026-10-19 (columns excerpted; the 1h cache-write column omitted).}
  \note{[2 min] [CORE] A real price sheet, excerpted from the provider's published pricing page as of today; one column is omitted for space. Read the first row: \$3 per million fresh input tokens, \$0.30 per million cached input tokens (cache hits), \$15 per million output tokens. That is exactly the tariff the chapter's notebooks use. Writing a prefix into the cache costs 1.25 times the base input price for a five-minute lifetime. The footnote matters for the next two slides: on the newer models a cache hit costs 5 percent of fresh input, not 10. ENGAGE: why does a provider charge less for a cached prefix? (Expected: the server keeps the prefix's attention state and does not recompute it; the discount applies only to an exact, resident prefix.) WARN: tariffs change; the arithmetic that follows takes the tariff as an input and must be rerun when it moves. FLEX: [CORE].}
\end{frame}

\begin{frame}{Prefix caching cuts the 30-turn bill 5.6$\times$; the quadratic term stays}
  \centering
  \mlsyssteps[width=0.92\textwidth,height=4.2cm,keepaspectratio]{resend_step}{2}\\[0.05cm]
  {\small Cached re-reads are still 50.8\% of the cached bill. If 10\% of the prefix misses, the bill rises 45.7\%.}
  \mlsyscredit{Agent Economics, Napkin Math 19.2: 8,000-token first prompt, 1,500 appended per turn, 200 output tokens per call (hypothetical). Values: images/calcs.py.}
  \note{[2.5 min] [CORE] The repair trajectory's token bill, call by call. Step 1, gray: an append-only transcript resent at the fresh price. Each call reads everything before it, so input grows every turn and the 30 calls read 892,500 tokens in all; the bill is \$2.767. Step 2, blue: with ideal prefix reuse only 51,500 tokens are fresh, the first prompt plus each turn's increment; the rest is billed at the cached rate. The bill is \$0.497, 5.6 times less. But the blue bars still rise: the cached term grows with the square of the horizon, at a tenth of the slope, and those re-reads are 50.8 percent of what remains. ENGAGE: a timestamp moves to the front of the system prompt. What happens? (Expected: the prefix breaks at its first changed token every call, and a 10 percent miss alone adds \$0.227, 45.7 percent.) WARN: caching changes the bill, not the tokens the model reads, and not acceptance. LINK: Lecture 2 owns the resend arithmetic; this slide prices it. FLEX: [CORE].}
\end{frame}

\begin{frame}{Compaction pays only if enough calls remain to repay it}
  \centering
  \safeimg[width=0.92\textwidth,height=4.5cm,keepaspectratio]{payback.pdf}\\[0.05cm]
  {\footnotesize Compaction replaces old turns with a summary.}\quad{\small $n^\star = 1 + \left\lceil \dfrac{p_{\text{out}}K_{\text{sum}} + p_{\text{in}}S_c}{p_{\text{cache}}(S - S_c)} \right\rceil$}
  \mlsyscredit{Agent Economics, \S19.3.2 and Napkin Math 19.3 (hypothetical; tariff of Napkin Math 19.2). Curves and payback: images/calcs.py.}
  \note{[3 min] [CORE] Compaction rewrites the transcript, so the next call misses the cache. Set the two paths side by side. Compacting costs a summary's output, 1,500 tokens at \$15 per million, plus resending the 12,000-token summary fresh: 5.85 cents up front. Every later call then reads 18,500 fewer cached tokens, saving 0.555 cents. The blue line is the difference, compact minus keep, against the calls that remain; it crosses zero at 12 calls. With 15 calls left, compaction saves 1.92 cents. Now apply the newer tariff's 5 percent cache rate: the saving per call halves, the orange line crosses at 23 calls, and with 15 left compaction loses 1.97 cents. ENGAGE: why does a cheaper cache push the profitable compaction point earlier? (Expected: the upfront cost is priced at fresh and output rates; only the saving is priced at the cache rate.) WARN: a summary that drops a constraint and lowers acceptance erases any token saving. LINK: Lecture 4 decides when to compact for context limits; this slide asks only whether it pays on price. FLEX: [CORE].}
\end{frame}

\begin{frame}{Estimate: retry cheap, or call strong once?}
  \mlsysactivity{Estimate}{3 min}
  \begin{columns}[T]
    \begin{column}{0.40\textwidth}
      \begin{mlsyscard}{mlsysaccent}
        \footnotesize\textbf{Given}\\[0.03cm]
        Cheap: $p = 0.30$, \$0.10, up to 5 tries\\
        Strong: $p = 0.70$, \$0.40, one try\\
        Attempts independent
      \end{mlsyscard}
      \vspace{0.1cm}
      {\footnotesize Cost per accepted task: retry until a pass, run 5 in parallel, or call the strong model once?}
    \end{column}
    \begin{column}{0.56\textwidth}
      \small
      \wstep{2}{Coverage $1 - 0.7^5 = 83.2\%$; mean attempts $0.832/0.3 = 2.773$}
      \wstep{2}{Sequential: $c/p = 0.10/0.30 =$ \textbf{\textcolor{figgreen}{\$0.333}} per accepted task}
      \wstep{3}{Parallel batch: $5 \times 0.10 / 0.832 =$ \textbf{\$0.601}, 1.8$\times$ more}
      \wstep{3}{Strong single call: $0.40/0.70 =$ \textbf{\$0.571}}
      \uncover<3->{\footnotesize\textcolor{midgray}{With early stopping, the cap changes coverage, not cost per accepted task.}}
    \end{column}
  \end{columns}
  \mlsyscredit{Agent Economics, \S19.3.3 and Napkin Math 19.4 (hypothetical). Values: images/calcs.py.}
  \note{[3 min] [CORE] Pairs, two minutes. Step 2: sequential retry with early stopping covers 83.2 percent of tasks with 2.773 attempts on average, \$0.277 per submitted task and \$0.333 per accepted task. The expected spend and the expected accepted count both scale with coverage, so their ratio is the attempt cost over the pass probability whatever the cap. Step 3: launching all five at once pays for all five every time, \$0.601 per accepted task, 1.8 times more, the ratio of 5 to 2.773 attempts. One strong call costs \$0.571 and covers only 70 percent. ENGAGE: when is the parallel burst worth 1.8 times the money? (Expected: when a deadline values the shorter wall-clock time.) WARN: the independence assumption is the weak point; the next slide shows real tasks violate it. A larger reasoning budget is priced the same way: it pays only if acceptance rises by a larger proportion than attempt cost. FLEX: [CORE].}
\end{frame}

\begin{frame}{Real coverage grows far slower than independent retries predict}
  \begin{columns}[c]
    \begin{column}{0.50\textwidth}
      \centering\safeimg[width=\textwidth,height=5.3cm,keepaspectratio]{coverage.pdf}
    \end{column}
    \begin{column}{0.47\textwidth}
      \small Repair tasks on SWE-bench Lite, 250 samples each.\\[0.2cm]
      One sample solves 16\% of tasks. The independent model predicts 94\% by 16 samples; the measured coverage is 39\%, and 56\% at 250.\\[0.2cm]
      {\footnotesize Hard tasks stay hard on every try, so an average $p$ overstates coverage and understates $\Ceff$.}
    \end{column}
  \end{columns}
  \mlsyscredit{Brown et al., ``Large Language Monkeys,'' 2024, recomputed from released per-sample labels (DeepSeek-Coder-V2-Instruct with Moatless Tools, all 300 tasks), via MLSysIM. Prediction: images/calcs.py.}
  \note{[1.5 min] [OPTIONAL] Real data. Brown and colleagues sampled 250 repair attempts per task on SWE-bench Lite and released which samples passed. Violet: the fraction of tasks solved by at least one of k samples. Gray dashes: what independent attempts with the measured single-sample rate, 16 percent, would predict. At 16 samples the model says 94 percent; the measurement says 39 percent. Coverage keeps climbing roughly with the logarithm of samples instead of saturating. ENGAGE: what does this do to the previous slide's \$0.333? (Expected: on real task mixes, retries buy less coverage than predicted, so realized cost per accepted task drifts above attempt cost over pass rate.) WARN: more samples also raise the chance that a flawed candidate passes the tests, which Lecture 3 priced as verifier precision. Land the segment: reuse lowers a coefficient and pays back over a horizon; retries pay by acceptance per dollar, measured per task. FLEX: [OPTIONAL].}
\end{frame}

% =============================================================================
% SECTION: SPEED AND ROUTING
% =============================================================================
\mlsysnote{[0.5 min] [CORE] Reuse and retries decide what the runtime buys. Next: where faster execution shortens a task, and which model should attempt it at all.}
\mlsysdivider{Where does faster inference pay, and which model should run?}{Agent Economics, \S19.4 and \S19.5}
\section{Speed}

\begin{frame}{Predict: what does a 4$\times$ faster model buy?}
  \mlsysactivity{Predict}{2 min}
  \begin{columns}[c]
    \begin{column}{0.79\textwidth}
      \centering
      \mlsyssteps[width=\textwidth,height=3.9cm,keepaspectratio]{waterfall_step}{4}\\[0.05cm]
      {\small Model calls take 29.25 s of a 312.40 s episode. Episode speedup?}\\[0.08cm]
      {\normalsize\choice{A}{4$\times$}\quad\choice{B}{about 2$\times$}\quad\answerat{2}{C}{about 1.1$\times$}}
    \end{column}
    \begin{column}{0.19\textwidth}
      \centering\safeimg[width=\textwidth,height=2.7cm,keepaspectratio]{gene-amdahl.jpg}\\[0.05cm]
      {\footnotesize Gene Amdahl, 2008}\\[0.15cm]
      \uncover<2->{\footnotesize Critical-path acceleration ceiling: $1/(1-f) = 1.10\times$ at $f = 9.4\%$}
    \end{column}
  \end{columns}
  \mlsyscredit{Agent Economics, \S19.4.2, Napkin Math 19.5 and its waterfall figure (hypothetical episode), redrawn. Photo: Perry Kivolowitz, CC BY 3.0, Wikimedia Commons.}
  \note{[3 min] [CORE] A short serial repair episode measured as disjoint intervals: 245.60 s of required tests, 37.35 s of sandbox creation and snapshot export, 0.20 s of runtime setup, and 29.25 s of model calls. ENGAGE: vote. Step 2: C. Model calls are 9.4 percent of the episode. Amdahl's law: the new time is the old time times one minus that share plus the share over the speedup. A 4 times faster model gives 290.46 s, a saving of 21.94 s, 1.08 times. Step 3: speculative decoding at 3.29 times on the decode stream gives 292.05 s, 1.07 times. Step 4: instant model calls stop at the red line, 283.15 s, 1.10 times; the ceiling is one over one minus 9.4 percent. Halving the test time instead saves 122.8 s and lowers the sandbox lease charge too. WARN: the ceiling needs only a baseline profile, so compute it before buying faster serving. LINK: Lecture 18 taught Amdahl's law for delegation; Lecture 14 measures the critical path from recorded spans. On a branching trace, every speedup is priced against the recomputed critical path, not by subtracting component savings. FLEX: [CORE].}
\end{frame}

\begin{frame}{Speculative decoding pays only when accepted drafts beat drafting}
  \begin{columns}[c]
    \begin{column}{0.52\textwidth}
      \centering
      {\Large $S_{\text{spec}} = \dfrac{m+1}{\gamma r + 1}$}\\[0.3cm]
      {\footnotesize $\gamma$: tokens drafted per round\quad $r$: draft cost per token, as a share of a target step\quad $m$: draft tokens accepted per round}\\[0.25cm]
      {\small Faster only if $m > \gamma r$}
    \end{column}
    \begin{column}{0.45\textwidth}
      \small $\gamma = 5$, $r = 0.08$: break-even at $m = 0.4$\\[0.2cm]
      \textcolor{figgreen}{\textbf{$m = 3.6$:}} decode 3.29$\times$ faster; episode 292.05 s, 1.07$\times$\\[0.2cm]
      \textcolor{figred}{\textbf{$m = 0.3$:}} decode 0.93$\times$; episode 314.65 s, slower\\[0.2cm]
      {\footnotesize Two gates: the serving condition under real load, and a decode share large enough to matter.}
    \end{column}
  \end{columns}
  \mlsyscredit{Agent Economics, \S19.4.4 and equation 3 (hypothetical draft rates; Leviathan et al. 2023; Chen et al. 2023). Values: images/calcs.py.}
  \note{[1.5 min] [OPTIONAL] A small draft model proposes gamma tokens per round; the large target model checks them in one pass, keeps the longest matching run, and adds one token of its own. Each round yields m plus 1 tokens for the cost of gamma times r plus one target steps. With 5 drafted tokens at 0.08 of a step each, break-even is 0.4 accepted tokens per round. If drafts align and 3.6 are accepted, decode runs 3.29 times faster, but the whole episode only 1.07 times, because decode is a small share of it. If a hard reasoning task drops acceptance to 0.3, decode slows to 0.93 times and the episode grows to 314.65 s. ENGAGE: why might the serving gain vanish under heavy load? (Expected: at large batch, checking extra positions costs arithmetic the server no longer has spare.) WARN: a high draft acceptance rate is a serving metric, not a trajectory result. FLEX: [OPTIONAL].}
\end{frame}

\begin{frame}{Price is a poor guide to score on long software tasks}
  \centering
  \safeimg[width=0.92\textwidth,height=4.4cm,keepaspectratio]{frontierswe.pdf}\\[0.05cm]
  {\small 19 configurations: \$9 to \$1,030 per trial. The top score costs 10$\times$ the runner-up for 3 points; a median run-to-run spread is 21 points.}
  \mlsyscredit{FrontierSWE V2 leaderboard (frontierswe.com), 34 tasks, 5 trials each, 20-hour budget per trial; via Epoch AI Benchmarking Hub and MLSysIM. Ratios: images/calcs.py.}
  \note{[1.5 min] [OPTIONAL] Real data from a long-horizon software benchmark: every configuration runs each task five times under a 20-hour budget. Each dot is the mean score of five trials, the bar spans the worst to the best trial, and the axis is the average cost of one trial on a log scale. Costs span \$9 to \$1,030; trials take 1 to 19 hours. The highest-scoring configuration costs 10 times as much per trial as the runner-up for 3 more points, while the median gap between a configuration's best and worst trial is 21 points. ENGAGE: what does that spread mean for routing? (Expected: a tier's acceptance must be measured on the tasks that reach it, not read from its price or a headline score.) WARN: this is cost per trial, not per accepted task; dividing by acceptance would reorder some points. FLEX: [OPTIONAL].}
\end{frame}

\begin{frame}{A cascade saves money only through its early exits}
  \centering
  \mlsyssteps[width=0.95\textwidth,height=5.0cm,keepaspectratio]{cascade_step}{3}\\[0.02cm]
  {\small $\mathbb{E}[C] = \sum_j \rho_j (c_j + g_j)$,\quad $\rho_j$: share of tasks reaching tier $j$}
  \mlsyscredit{Agent Economics, \S19.5.1, Napkin Math 19.6 and its cascade figure (hypothetical, 100,000 code-editing tasks), redrawn. Values: images/calcs.py.}
  \note{[3 min] [CORE] A three-tier cascade on 100,000 code-editing tasks, every exit judged by the same gate. Step 1: tier 1 accepts 60,000 and passes 40,000 on; tier 2 accepts 28,000 of them and passes 12,000; the frontier tier accepts 11,400 and leaves 600. Conditional acceptance is 60, 70, and 95 percent, measured on the tasks each tier actually receives. Step 2: the cascade spends \$992 for 99,400 accepted, \$0.00998 each; sending everything to the frontier model costs \$6,100 for 97,000, \$0.06289, 6.3 times more. Nearly all the gap comes from the 60 percent that exit at tier 1. Step 3: the price is latency. A task that climbs every tier pays every gate in series, 6.75 s against 4.90 s direct, though the mean is 1.72 s. ENGAGE: how would you remove that penalty for urgent tasks? (Expected: route tasks predicted to fail early tiers straight to the frontier tier.) WARN: later tiers see a filtered, harder remainder, so their acceptance cannot be copied from full-workload benchmarks. FLEX: [CORE].}
\end{frame}

\begin{frame}{Diagnose: when does a cheaper gate cost more?}
  \mlsysactivity{Diagnose}{3 min}
  \begin{columns}[c]
    \begin{column}{0.56\textwidth}
      \small A team replaces tier 1's test gate with a lint check. The cascade log shows the same exits at lower cost.\\[0.2cm]
      Measured later: \textbf{0.9\%} of tier-1 exits break an independent check, and each costs \textbf{\$10} to fix in production.\\[0.25cm]
      {\normalsize Is the cascade still cheaper than direct frontier?}
    \end{column}
    \begin{column}{0.41\textwidth}
      \uncover<2->{\begin{mlsyscard}{figred}\small
      Harm: $60{,}000 \times 0.009 \times \$10 = \$5{,}400$\\[0.05cm]
      Cascade: $(\$992 + \$5{,}400) / 99{,}400 = \mathbf{\$0.06431}$\\[0.05cm]
      Direct frontier: \textbf{\$0.06289}\\[0.1cm]
      \footnotesize The \$0.053 saving per accepted result is gone.
      \end{mlsyscard}}
    \end{column}
  \end{columns}
  \mlsyscredit{Agent Economics, \S19.5.2 and \S19.10 (fallacy: a cheap first tier makes savings independent of its gate). Values: images/calcs.py.}
  \note{[3 min] [CORE] Pairs, two minutes. The log records a sub-cent exit; it cannot record that the exit was wrong. Step 2: tier 1 exits 60,000 tasks; 0.9 percent of them, 540, are false accepts at \$10 each, \$5,400 of harm. Added to the \$992 cascade spend over 99,400 accepted results, that is \$0.06431 each, above the direct frontier's \$0.06289. The harm term adds the probability of a false accept, over initiated tasks, times its harm. ENGAGE: a team measures false accepts only among accepted results. What must it multiply by to get the per-task probability? (Expected: the overall acceptance rate.) WARN: every exit must satisfy the same acceptance contract; a cheap tier never earns a weaker gate, and catastrophic harms stay under authority limits whatever their expected value. LINK: Lecture 3 owns verifier precision. Land the segment: a speedup pays only on the critical path, and a cascade pays only through early exits under one gate. FLEX: [CORE].}
\end{frame}

% =============================================================================
% SECTION: CAPACITY
% =============================================================================
\mlsysnote{[0.5 min] [CORE] So far, one task at a time. Now many trajectories share model servers, sandboxes, and reviewers.}
\mlsysdivider{How much capacity does a goodput target need?}{Agent Economics, \S19.6}
\section{Capacity}

\begin{frame}{One arrival rate makes different demands on each resource}
  \begin{columns}[c]
    \begin{column}{0.36\textwidth}
      \centering\mlsysphoto[width=\textwidth,height=3.6cm,keepaspectratio]{datacenter-server-racks.jpg}\\[0.05cm]
      {\footnotesize Model servers, sandboxes, and quotas are separate pools.}
    \end{column}
    \begin{column}{0.61\textwidth}
      \small 1 trajectory per minute, 20 minutes each, 60 calls of 30,000 tokens\\[0.15cm]
      \wstep{1}{Little's law: $L = \lambda W = 1 \times 20 =$ \textbf{20} trajectories present}
      \wstep{1}{Calls: \textbf{60} per minute; tokens: \textbf{1.8M} per minute, \textbf{90\%} of a 2M quota}
      \wstep{2}{Goodput $\mathcal{G}$ = accepted tasks per unit time; admit $\lambda = \mathcal{G}/\alpha$}
      \wstep{2}{Acceptance falls 10\%: demand rises to \textbf{\textcolor{figred}{100\%}} of the quota}
    \end{column}
  \end{columns}
  \mlsyscredit{Agent Economics, \S19.6.1 and Napkin Math 19.7 (hypothetical fleet). Values: images/calcs.py. Photo: Carl Lender, CC BY 2.0, Wikimedia Commons.}
  \note{[3 min] [CORE] Admitting one 20-minute repair trajectory a minute. Step 1: Little's law gives 20 trajectories present on average, each holding a sandbox and retained context. They send 60 model calls a minute. At 30,000 quota-counted tokens per call, demand is 1.8 million tokens a minute, 90 percent of a 2 million allowance; the quota alone caps admission at 1.11 per minute. Apply Little's law at each boundary separately: a trajectory waiting five minutes on a test uses no accelerator but still holds its sandbox. Step 2: trajectory goodput, accepted tasks per unit time, is what the platform delivers. Holding goodput fixed, admissions scale as one over acceptance, so a 10 percent relative drop in acceptance pushes demand to the full quota with no room for bursts. ENGAGE: a prompt change lowers acceptance but nobody touches the fleet. What breaks? (Expected: offered load rises; the quota or the server pool saturates.) WARN: a model or prompt change that lowers acceptance is also a capacity change. LINK: Lecture 5 owns Little's law. FLEX: [CORE].}
\end{frame}

\begin{frame}{Quota rules differ by contract, so read the actual one}
  \begin{columns}[c]
    \begin{column}{0.60\textwidth}
      \begin{logblock}[7.2]
        Many API providers use a combined "tokens per minute" (TPM) limit that may include all tokens, both cached and uncached, input and output. \textbf{For most Claude models, only uncached input tokens count toward your ITPM rate limits.}\\[0.12cm]
        * `cache\_read\_input\_tokens` (tokens read from cache) \textbf{Do NOT count toward ITPM} for most models
      \end{logblock}
    \end{column}
    \begin{column}{0.37\textwidth}
      \small The previous slide assumed a quota that counts cached input at full weight.\\[0.2cm]
      Under this contract, caching also raises the admission ceiling.
    \end{column}
  \end{columns}
  \mlsyscredit{Anthropic, ``Rate limits,'' platform.claude.com/docs/en/api/rate-limits, \S Cache-aware ITPM, retrieved 2026-10-19 (excerpt; the check mark is dropped). ITPM: input tokens per minute.}
  \note{[1 min] [OPTIONAL] A real rate-limit page, excerpted. This provider counts only uncached input toward its input-token-per-minute limit, while other providers count all tokens. The chapter's notebook assumed a quota that counts cached input at full weight, and under that contract prefix caching lowers the bill without raising the admission ceiling. Under this one, caching does both. ENGAGE: which number on the previous slide changes under this contract? (Expected: the 90 percent quota use falls, because most of each call's input is a cached prefix.) WARN: limits are counted over windows of different lengths and in requests, tokens, or connections; admission control must match the contract's semantics, measured against production traces. FLEX: [OPTIONAL].}
\end{frame}

\begin{frame}{Filling a server raises throughput and slows every user}
  \begin{columns}[c]
    \begin{column}{0.50\textwidth}
      \centering\safeimg[width=\textwidth,height=5.3cm,keepaspectratio]{inferencex.pdf}
    \end{column}
    \begin{column}{0.47\textwidth}
      \small Measured: one 8-GPU server, one model, concurrency 4 to 64.\\[0.2cm]
      16$\times$ the users gives \textbf{5.1$\times$} the tokens per GPU, and each user's token time rises from \textbf{9.8 ms} to \textbf{31.5 ms}.\\[0.2cm]
      {\footnotesize A trajectory pays that wait on every serial call.}
    \end{column}
  \end{columns}
  \mlsyscredit{SemiAnalysis InferenceX (github.com/SemiAnalysisAI/InferenceX), DeepSeek-R1 FP8, 1,024 in and 1,024 out, 8$\times$ H200, SGLang, run of 2026-05-20, via MLSysIM. Ratios: images/calcs.py.}
  \note{[1.5 min] [OPTIONAL] Real measurements from an open serving benchmark: one model on one eight-GPU server, swept from 4 to 64 concurrent users. Across: the tokens per second each user sees, falling to the right. Up: tokens per second per GPU, which is what the operator pays for. Sixteen times the users buys 5.1 times the throughput per GPU, and each user's time per token rises from 9.8 ms to 31.5 ms. ENGAGE: an operator sees the GPU bill drop as the server fills. What does the agent team see? (Expected: every one of a trajectory's serial calls takes longer, while sandboxes and context stay held.) WARN: this is batching on one server, not a queue; the next slide adds the queue wait in front of it. Both push the same way. FLEX: [OPTIONAL].}
\end{frame}

\begin{frame}{Predict: from 70\% to 90\% busy, how much wait?}
  \mlsysactivity{Predict}{3 min}
  \begin{columns}[c]
    \begin{column}{0.38\textwidth}
      \centering\safeimg[width=\textwidth,height=2.6cm,keepaspectratio]{checkout-queue.jpg}\\[0.1cm]
      {\footnotesize $\mathbb{E}[W] \approx \dfrac{c_a^2 + c_s^2}{2}\,\dfrac{\rho}{1-\rho}\,\mathbb{E}[S]$}\\[0.05cm]
      {\scriptsize $c^2$: burstiness of arrivals and service; $\rho$: utilization; $S$: service time 1 s}
    \end{column}
    \begin{column}{0.59\textwidth}
      \centering
      \mlsyssteps[width=\textwidth,height=3.9cm,keepaspectratio]{kingman_step}{2}\\[0.03cm]
      {\small\choice{A}{1.3$\times$}\quad\choice{B}{2$\times$}\quad\answerat{2}{C}{about 4$\times$}}\\[0.03cm]
      \uncover<2->{\footnotesize 28.6\% more throughput costs 3.9$\times$ the wait.}
    \end{column}
  \end{columns}
  \mlsyscredit{Agent Economics, \S19.6.2, equation 5 (Kingman 1961) and Napkin Math 19.8 ($c_a^2 = 2$, $c_s^2 = 4$, hypothetical). Values: images/calcs.py. Photo: Gerald England, CC BY-SA 2.0, Wikimedia Commons.}
  \note{[3 min] [CORE] A checkout queue is the picture: when the cashier is busy most of the time, a small rise in shoppers makes the line long. Kingman's approximation for one server: the mean wait is the burstiness term times rho over one minus rho times the service time. With burstiness 3 and 1 s of service, ENGAGE: vote. Step 2: C. At 70 percent the wait is 7 s; at 90 percent, 27 s, 3.9 times longer, for 28.6 percent more throughput. At \$0.002 per second of held resources, the holding charge per call rises from \$0.014 to \$0.054. A trajectory pays it on every serial step. Tails compound too: if each of 30 serial calls exceeds its 99th percentile with probability 0.01, the trajectory meets at least one such delay with probability 0.26. LINK: this is Lecture 1's horizon compounding applied to latency. WARN: the formula is a heavy-traffic estimate of the mean for one server; a shared pool needs a pool model or a trace-driven simulation. Halving burstiness halves the wait at any utilization, a cheaper lever than buying capacity. FLEX: [CORE].}
\end{frame}

\begin{frame}{Owning capacity wins only above the break-even load}
  \begin{columns}[c]
    \begin{column}{0.50\textwidth}
      \centering\safeimg[width=\textwidth,height=5.3cm,keepaspectratio]{breakeven.pdf}
    \end{column}
    \begin{column}{0.47\textwidth}
      \small $u^\star = \dfrac{F}{3600\,R\,(P_{\text{eff}} - v)}$\\[0.15cm]
      {\footnotesize $F$: \$28 per hour fixed;\ $R$: 1,200 output tokens/s;\ $v$: \$2 per M output;\ $P_{\text{eff}}$: API price per output token with its input}\\[0.2cm]
      80\% cached input: $P_{\text{eff}} = \$19.2$/M, $u^\star = 37.7\%$; at 52\%, owning saves \$10.64 an hour\\[0.1cm]
      No caching: $u^\star = 23.1\%$
    \end{column}
  \end{columns}
  \mlsyscredit{Agent Economics, \S19.6.3, equation 6 and Napkin Math 19.9 (hypothetical; Armbrust et al. 2010). Values: images/calcs.py.}
  \note{[1.5 min] [OPTIONAL] Rent or own. The API line starts at zero and rises with load; the dedicated line starts at its fixed \$28 an hour and rises at the \$2 per million marginal cost. With 5 input tokens per output token and 80 percent of input cached, the API's effective price is \$19.2 per million output tokens. The lines cross at 37.7 percent of the reference rate. At 52 percent load, 2,246,400 output tokens an hour cost \$43.13 through the API and \$32.49 owned. Without caching the API's effective price is higher and owning pays from 23.1 percent. ENGAGE: why does a provider's cache discount weaken the case for owning? (Expected: it lowers the API price per output token, which moves the crossing right.) WARN: the comparison needs matched capability, privacy, and tail latency, and busy-hour load is not around-the-clock load; idle hours pay the fixed cost. Land the segment: capacity is sized for goodput at its acceptance, and utilization is paid for in waiting. FLEX: [OPTIONAL].}
\end{frame}

% =============================================================================
% SECTION: BUDGETS AND DESIGNS
% =============================================================================
\mlsysnote{[0.5 min] [CORE] Capacity bounds the fleet. Now bound what one task and its children may spend, and then price whole architectures.}
\mlsysdivider{How does a spending limit hold, and which design pays?}{Agent Economics, \S19.7 to \S19.9}
\section{Budgets}

\begin{frame}[fragile]{A limit checked after the charge cannot stop the charge}
  \begin{columns}[c]
    \begin{column}{0.60\textwidth}
      \begin{logblock}[6.9]
        per\_instance\_cost\_limit: float = Field(\\
        \ind{4}default=3.0,\\
        ...\\
        def \_update\_stats(self, *, input\_tokens: int,\\
        \ind{17}output\_tokens: int, cost: float) -> None:\\
        \ind{4}...\\
        \ind{4}\textbf{self.stats.instance\_cost += cost}\\
        \ind{4}...\\
        \ind{4}if 0 < self.config.per\_instance\_cost\_limit < self.stats.instance\_cost:\\
        \ind{8}...\\
        \ind{8}msg = "Instance cost limit exceeded"\\
        \ind{8}raise InstanceCostLimitExceededError(msg)
      \end{logblock}
    \end{column}
    \begin{column}{0.37\textwidth}
      \small A 2024 run of this agent capped each task at \$4. Of its 500 tasks, \textbf{90} ended above \$4, the highest at \textbf{\$4.20}.\\[0.2cm]
      \uncover<2->{\begin{mlsyscard}{mlsysaccent}\small Reserve a bound $\bar c$ before dispatch:\\ settled $+$ reserved $\le$ budget\end{mlsyscard}}
    \end{column}
  \end{columns}
  \mlsyscredit{SWE-agent, sweagent/agent/models.py, commit 3ea751c (excerpt). \$4 cap: Kapoor et al., ``AI Agents That Matter,'' 2024, p.~1. Run A costs via MLSysIM; counts: images/calcs.py.}
  \note{[2.5 min] [CORE] Real code from an open-source repair agent, excerpted; the dots mark omitted lines. Read it: after each model call returns, the agent adds the call's cost to the task's total, and only then compares the total with the limit. The check stops the next call, not this one. The 2024 run on the earlier slide used this agent with a \$4 cap per task, reported by Kapoor and colleagues; 90 of its 500 trajectories ended above \$4, the highest at \$4.20. The cap held approximately because single calls were small. Step 2: the runtime's invariant. Before dispatch, reserve an enforceable bound on the call's maximum charge, including its maximum completion and hidden reasoning tokens; settled plus reserved never exceeds the budget; settlement releases the unused part. ENGAGE: a provider's cache discount is likely but not guaranteed. What bound do you reserve? (Expected: the undiscounted price.) LINK: Lecture 10 owns bounded-liability admission for one trajectory; this segment extends it across children. FLEX: [CORE].}
\end{frame}

\begin{frame}{Predict: when may the root release the money?}
  \mlsysactivity{Predict}{3 min}
  \centering
  \mlsyssteps[width=0.92\textwidth,height=4.5cm,keepaspectratio]{budget_step}{5}\\[0.05cm]
  \only<1-3|handout:0>{\small Fuzzing child at its stop: \$0.85 settled, \$1 reserved for a remote job.}%
  \only<4->{\small A release at local stop would return \$1.65; only \textbf{\$0.65} is free.}
  \mlsyscredit{Agent Economics, \S19.7.2, table 5, Napkin Math 19.10 (hypothetical) and its conservation figure, redrawn. Values: images/calcs.py.}
  \note{[3.5 min] [CORE] One repair task with a \$5 budget; each bar is the root's account and always sums to \$5. Step 1: planning settles \$0.042. Step 2: the root grants \$1.20 to a trace-analysis child and \$2.50 to a fuzzing child; both move from free to reserved. Step 3: the trace child closes with \$0.38 settled; its unused \$0.82 returns to free. ENGAGE at step 3, pairs: the fuzzing child has settled \$0.85 and holds \$1 for a remote operation when its local process stops. How much may the root release now? Step 4: nothing moves. Stopping closes admission but releases no liability; the whole \$2.50 stays reserved. A release at local stop would have returned \$1.65; only \$0.65 was truly free. Step 5: the remote job bills \$0.70 late; the child closes with \$1.55 settled and \$0.95 returns. Had the root released \$1.65, the late bill would have landed on money promised elsewhere. WARN: redelivered bills must settle once, keyed by operation identity. LINK: Lecture 18 owns subtree cancellation; Lecture 7 owns settlement before retry. FLEX: [CORE].}
\end{frame}

\begin{frame}{Price designs by cost per resolved ticket, human fallback included}
  \begin{columns}[c]
    \begin{column}{0.64\textwidth}
      \centering
      \mlsyssteps[width=\textwidth,height=3.6cm,keepaspectratio]{arch_step}{2}\\[0.05cm]
      {\footnotesize 10,000 upgrades a month; rejects go to a \$35 engineer who resolves 90\%. The loop accepts 81.73\% in 1.711 turns.}
    \end{column}
    \begin{column}{0.33\textwidth}
      {\footnotesize\textbf{The workload, in real life}}\\[0.05cm]
      \begin{logblock}[7.2]
        Bump actions/checkout from 4.1.4 to 4.1.6 in the github-actions group\\[0.05cm]
        dependabot[bot] opened 2024-06-01 15:55:46Z, merged 16:01:47Z\\
        +3 $-$3, 2 files
      \end{logblock}
      \vspace{0.1cm}
      {\scriptsize A deterministic rewriter, pallets/flask \#5493}
    \end{column}
  \end{columns}
  \mlsyscredit{Agent Economics, \S19.8.3, table 7 and Napkin Math 19.11 (hypothetical). Pull request: github.com/pallets/flask/pull/5493 (GitHub API). Values: images/calcs.py.}
  \note{[3 min] [CORE] Dependency upgrades are a real workload; the excerpt is a real pull request opened and merged by a deterministic rule-based rewriter. The notebook compares three designs under one contract: a rule-based rewriter that passes 32 percent, one frontier call that passes 68 percent, and a three-turn loop on a small model whose turns pass 55, 42, and 30 percent of what reaches them, 81.73 percent in all. Step 1: automated spend per automated acceptance ranks the rewriter best, \$0.00937. Step 2: every rejected ticket goes to a \$35 engineer. Per resolved ticket the rewriter costs \$25.54 and misses the 95 percent service floor; the single call costs \$11.62; the loop costs \$6.53. Each point of automated acceptance saves \$0.35 of fallback per ticket, more than any design's automated spend. ENGAGE: when does adding a turn pay, if the engineer resolves everything? (Expected: when the turn's conditional acceptance times \$35 exceeds the turn's cost.) WARN: a ticket routed to a person counts only when the person finishes it under the same tests. A one-time adaptation is priced the same way: the chapter's smaller model repays \$12,500 after 56,510 tasks. LINK: Lecture 1 owns workflow against loop; this prices the choice. FLEX: [CORE].}
\end{frame}

% =============================================================================
% SYNTHESIS AND CLOSE
% =============================================================================
\begin{frame}{Back to the hook: each workload is bound by a different term}
  \centering\scriptsize
  \begin{tabular}{@{}lrrrr@{}}
    \toprule
    & \textbf{Repair} & \textbf{Research} & \textbf{Remediation} & \textbf{Support}\\
    \midrule
    Attempt cost & \$1.50 & \$1.05 & \$20.45 & \$0.65\\
    Acceptance & 0.40 & 0.60 & 0.70 & 0.85\\
    $\Ceff$ & \$3.74 & \$1.75 & \$29.21 & \$0.76\\
    Largest term & sandbox and tools & fan-out tokens & human approval & escalation handoff\\
    \midrule
    Saving: halve token price & 16.6\% & \textbf{42.9\%} & 0.4\% & 2.3\%\\
    Saving: halve largest term & \textbf{30.1\%} & \textbf{42.9\%} & \textbf{48.9\%} & 46.2\%\\
    Saving: 10 more points accepted & 20\% & 14.3\% & 12.5\% & \textbf{65.6\%}\\
    \bottomrule
  \end{tabular}\\[0.25cm]
  {\small The hook's discount acted on one term. Acceptance divides them all.}
  \mlsyscredit{Agent Economics, Lighthouse 19.1, table 8 (scenario values; research fan-out from Hadfield et al. 2025). Savings recomputed: images/calcs.py.}
  \note{[2.5 min] [CORE] The four lighthouse workloads priced in one account. Repair: model tokens are 33.2 percent of the attempt; holding the sandbox and tools across 30 turns is the largest term. Halving the token price saves 16.6 percent; ten more points of acceptance save 20 percent. Research is the only token-bound workload, 85.7 percent of its attempt, so caching and bounded fan-out are its levers. Remediation is dominated by the engineer who approves the production change; token prices barely matter. Support is dominated by the \$4 handoff on each rejected session, so acceptance both raises the denominator and removes handoffs, 65.6 percent. Answer the hook: the low-price model cut the term that was not binding and lowered acceptance, which divides every term. ENGAGE: which workload would you move to a cheaper model first? (Expected: research, and only if acceptance holds.) Rule from the chapter: higher acceptance beats a halved token price whenever the model's share of the attempt is below twice the gain over the new acceptance, 40 percent for repair. LINK: the four lighthouse workloads come from Lecture 1. FLEX: [CORE].}
\end{frame}

\begin{frame}{Takeaways}
  \thumbline{thumb_swe_step2.pdf}{Each accepted task carries $(1-\alpha)/\alpha$ failures; a real run cost \$5.10 per resolved task against a \$1.16 survivor average.}\\[0.12cm]
  \thumbline{thumb_payback.pdf}{Caching cuts the 30-turn bill 5.6$\times$ but keeps the quadratic term; compaction pays after 12 calls at a 10\% cache rate, 23 at 5\%.}\\[0.12cm]
  \thumbline{thumb_waterfall_step4.pdf}{At 9.4\% model time, no model speedup passes 1.10$\times$. A cascade saves 6.3$\times$ only through early exits under one gate.}\\[0.12cm]
  \thumbline{thumb_kingman_step2.pdf}{Admit goodput over $\alpha$. From 70\% to 90\% utilization, 28.6\% more throughput costs 3.9$\times$ the wait.}\\[0.12cm]
  \thumbline{thumb_budget_step5.pdf}{Reserve before dispatch, release only at settlement, and price designs per resolved ticket with fallback.}
  \note{[1.5 min] [CORE] One takeaway per segment; read them as claims with their numbers. Ask students to underline the one that changes how they would read a model's price sheet. FLEX: [CORE].}
\end{frame}

\begin{frame}{Open questions}
  \begin{columns}[c]
    \begin{column}{0.48\textwidth}
      \small
      \begin{itemize}
        \item Benchmarks rank agents by accuracy. Should they rank by cost per accepted task, and under whose prices?
        \item Token prices fall 9$\times$ to 860$\times$ a year. Which term binds an agent in five years?
        \item Human fallback dominates many accounts. How should a team price a reviewer's queue and errors?
      \end{itemize}
    \end{column}
    \begin{column}{0.49\textwidth}
      \centering\safeimg[width=\textwidth,height=4.9cm,keepaspectratio]{kapoor.pdf}
    \end{column}
  \end{columns}
  \mlsyscredit{Data: Kapoor et al., ``AI Agents That Matter,'' arXiv:2407.01502, 2024, Table A1 (HumanEval, mean of five runs, April 2024 prices), redrawn.}
  \note{[1 min] [OPTIONAL] Real data, redrawn from the paper's table: Kapoor and colleagues reran coding agents five times on HumanEval and reported mean accuracy against mean dollar cost. Simple baselines (warming, retry, and escalation from a cheap model to expensive ones) match or beat complex agents at a fraction of the cost. Escalation, a cascade, reaches 85.0 percent for \$0.27; the debugging agent LDB on the same small model costs 2.3 times as much and scores 80.2 percent. The tree-search agent LATS on the large model costs 55 times as much as warming on the same model and scores lower. Discussion, two minutes with a neighbor on one of the three questions. For the first, point out that cost depends on prices that change, so a leaderboard would need to publish tokens, turns, and acceptance and let readers reprice. FLEX: [OPTIONAL].}
\end{frame}

\begin{frame}{Read and explore}
  \small
  \begin{itemize}
    \item \textbf{Agent Economics}, \S19.2 to \S19.9; Checkpoints 19.1 and 19.2; exercises 1 to 9
    \item Kapoor et al., \emph{AI Agents That Matter}, 2024: report cost beside accuracy
    \item Kingman, \emph{The Single Server Queue in Heavy Traffic}, 1961: three pages on why waiting explodes
    \item Dean and Barroso, \emph{The Tail at Scale}, 2013: tails dominate fan-out
    \item O'Neil, \emph{The Escrow Transactional Method}, 1986: concurrent reservations on one balance
    \item Armbrust et al., \emph{A View of Cloud Computing}, 2010: renting against owning
  \end{itemize}
  \note{[0.5 min] [CORE] Assign sections 19.2 to 19.9 and the two checkpoints. Exercise 3 (how many turns before a person takes over) and exercise 9 (a remediation budget under cancellation) practice the two decisions students find hardest. FLEX: [CORE].}
\end{frame}

\begin{frame}{What comes next?}
  \begin{columns}[c]
    \begin{column}{0.34\textwidth}
      \centering\safeimg[width=\textwidth,height=3.6cm,keepaspectratio]{thumb_next_stack.pdf}
    \end{column}
    \begin{column}{0.62\textwidth}
      {\small\textcolor{lightgray}{NEXT LECTURE \textbar{} CONCLUSION}}\\[0.5cm]
      {\LARGE\bfseries Who owns each transition of one trajectory, and what evidence permits the next?}
    \end{column}
  \end{columns}
  \note{[0.5 min] [CORE] Today priced every lever on one account. Lecture 20, the Conclusion, follows one remediation incident through every boundary of the system, names the owner and the record of each transition, and puts that attempt on today's account. Components that pass their own tests can still leave a task unfinished. FLEX: [CORE].}
\end{frame}

\end{document}
